\documentclass[11pt,letterpaper]{article}
\usepackage[utf8]{inputenc}

%----- Configuración del estilo del documento------%
\usepackage{epsfig,graphicx}
\usepackage[left=2cm,right=2cm,top=1.8cm,bottom=2.3cm]{geometry}
\usepackage{fancyhdr}
\usepackage{lastpage}
\usepackage{url}
\usepackage{tikz}
\pagestyle{fancy}
\fancyhf{}
\fancyhead[R]{\ifnum\value{page}>1 \textbf{Problema 8}\fi}
\rfoot{\textit{Página \thepage}}
%Encabezado izquierdo: logo
\lfoot{\includegraphics[height=1.0cm]{images.png}}
%------ Paquetes matemáticos básicos --------%
\usepackage{amsmath}
\usepackage{amssymb}
\usepackage{amsthm}

\usepackage[spanish]{babel}
\usepackage{graphicx}
\usepackage{hyperref}

\usepackage{tabularx}
\usepackage{xcolor}
\usepackage[table]{xcolor}
\usepackage{colortbl}
\usepackage{array, multirow, multicol, tabularx}
\usepackage{tcolorbox}
\newtheorem{theorem}{Theorem}[section]
\newtheorem{corollary}{Corollary}[theorem]
\newtheorem{lemma}[theorem]{Lema}

%------si-------%
\definecolor{B}{HTML}{FFFFFF}
\definecolor{G}{HTML}{5e5e5e}
\definecolor{R2}{HTML}{d53d40}
\definecolor{A2}{HTML}{034190}
\definecolor{V2}{HTML}{7faa50}
\newcommand{\R}{\mathbb{R}}
\newcommand{\C}{\mathcal{C}}
\newcommand{\N}{\mathbb{N}}
\newcommand{\Z}{\mathbb{Z}}
\newcommand{\Q}{\mathbb{Q}}
\renewcommand{\theenumi}{\Roman{enumi}}
\renewcommand{\labelenumi}{{\theenumi}.}
\definecolor{CIMAT}{HTML}{64293e}
\definecolor{R2}{HTML}{d53d40}
\definecolor{G}{HTML}{5e5e5e}
\newcommand{\dd}{\,dx}
\newcommand{\dt}{\,dt}
\newcommand{\du}{\,du}
\definecolor{CIMAT}{HTML}{64293e}
\definecolor{R2}{HTML}{d53d40}
\definecolor{G}{HTML}{5e5e5e}

\tcbset{
    colback=white,
    colframe=CIMAT,
    fonttitle=\bfseries,
    coltitle=white,
    arc=1mm,
    boxrule=0.6pt
}

\begin{document}
%---Evironment--%
	\makeatletter
        \renewenvironment{proof}[1][\proofname]{\par
            \pushQED{\qed}%
            \normalfont \topsep6\p@\@plus6\p@\relax
            \trivlist
            \item\relax
            {\itshape
            #1\@addpunct{.}}\par\vspace{\baselineskip}\ignorespaces
            }{%
            \popQED\endtrivlist\@endpefalse
            }
        \newenvironment{Solucion}[1][Soluci\'on]{\par
            \pushQED{\qed}%
            \normalfont \topsep6\p@\@plus6\p@\relax
            \trivlist
            \item\relax
            {\itshape
            #1\@addpunct{.}}\par\vspace{\baselineskip}\ignorespaces
            }{%
            \popQED\endtrivlist\@endpefalse
            }
            
    \makeatother

%------ Encabezado -------- %

\begin{center}
    \begin{minipage}{3cm}
    	\begin{center}
    		\includegraphics[height=3.4cm]{images.png}
    	\end{center}
    \end{minipage}\hfill
    	\begin{center}
    	\textbf{\large Taller de C\'alculo CIMAT}\\[0.1cm]
        \textbf{Facultad de Ciencias, ENES Juriquilla}\\[0.1cm]
        \textbf{Problema 8}\\[0.1cm]
         El\'ias L\'opez Rivera\\[0.1cm]
         Oscar Cach\'u Vel\'azquez \\[0.1cm]
         Francisco Javier Trejo Sanaphre \\[0.1cm]
    	\end{center}
\end{center}
\rule{17cm}{0.1mm}
%------ Fin de encabezado -------- %
\begin{tcolorbox}[title=\textcolor{white}{Problema 1}]
\textit{Muestre que los coeficientes de Fourier de $f$ existen, suponiendo que
$f:\R\to\R$ es Riemann integrable en $[-\pi,\pi]$.}
\end{tcolorbox}

\textbf{Demostración.}\,\
Para cada $n\in\Z$, el coeficiente complejo de Fourier de $f$ está definido por
\begin{equation*}
    c_n=\frac{1}{2\pi}\int_{-\pi}^{\pi} f(x)e^{-inx}.
\end{equation*}
Por la fórmula de Euler,
\begin{equation*}
    e^{-inx}=\cos(nx)-i\sen(nx).
\end{equation*}
Entonces
\begin{align*}
    c_n
    &=\frac{1}{2\pi}\int_{-\pi}^{\pi} f(x)\bigl(\cos(nx)-i\sen(nx)\bigr) \\
    &=\frac{1}{2\pi}\left(\int_{-\pi}^{\pi} f(x)\cos(nx)
    -i\int_{-\pi}^{\pi} f(x)\sen(nx)\right).
\end{align*}
Como $\cos(nx)$ y $\sen(nx)$ son funciones continuas en $[-\pi,\pi]$, y $f$ es Riemann integrable, los productos
\begin{equation*}
    f(x)\cos(nx)\qquad \text{y}\qquad f(x)\sen(nx)
\end{equation*}
también son Riemann integrables. Por tanto ambas integrales existen, y así $c_n$ existe para todo $n\in\Z$.
\hfill $\square$

\vspace{0.4cm}

\begin{tcolorbox}[title=\textcolor{white}{Problema 2}]
\textit{Muestre que, si $f$ es $2\pi$-periódica e integrable y $J$ es un intervalo de longitud $2\pi$, entonces}
\begin{equation*}
    \int_J f=\int_{-\pi}^{\pi}f.
\end{equation*}
\end{tcolorbox}
\textbf{Demostración.}\,\
Sea $J=[a,a+2\pi]$. Como $f$ es $2\pi$-periódica, se cumple
\begin{equation*}
    f(x+2\pi k)=f(x)\qquad \text{para todo }k\in\Z.
\end{equation*}
Tomemos $k\in\Z$ de tal forma que
\begin{equation*}
    c=a-2\pi k\in[0,2\pi).
\end{equation*}
Entonces $a=c+2\pi k$. Haciendo el cambio de variable $u=x-2\pi k$, obtenemos
\begin{equation*}
    \int_a^{a+2\pi} f(x)
    =\int_c^{c+2\pi} f(u+2\pi k)\,du
    =\int_c^{c+2\pi} f(u)\,du.
\end{equation*}
Ahora partimos la integral:
\begin{equation*}
    \int_c^{c+2\pi} f(u)\,du
    =\int_c^{2\pi} f(u)\,du+\int_{2\pi}^{c+2\pi} f(u)\,du.
\end{equation*}
En la segunda integral usamos el cambio $v=u-2\pi$. Como $f(v+2\pi)=f(v)$, queda
\begin{equation*}
    \int_{2\pi}^{c+2\pi} f(u)\,du
    =\int_0^c f(v+2\pi)\,dv
    =\int_0^c f(v)\,dv.
\end{equation*}
Por tanto
\begin{equation*}
    \int_a^{a+2\pi} f
    =\int_c^{2\pi}f+\int_0^c f
    =\int_0^{2\pi} f.
\end{equation*}
Finalmente, por el mismo argumento de periodicidad,
\begin{equation*}
    \int_{-\pi}^{\pi} f
    =\int_0^{2\pi} f.
\end{equation*}
Luego
\begin{equation*}
    \int_J f=\int_a^{a+2\pi}f=\int_{-\pi}^{\pi}f.
\end{equation*}
\hfill $\square$

\newpage

\begin{tcolorbox}[title=\textcolor{white}{Problema 3}]
\textit{Pruebe las siguientes afirmaciones:}
\begin{enumerate}
    \item[\textbf{a)}] Para $n\in\Z$,
    \begin{equation*}
        \frac{1}{2\pi}\int_{-\pi}^{\pi} e^{inx}=
        \begin{cases}
        0, & n\neq 0,\\
        1, & n=0.
        \end{cases}
    \end{equation*}
    \item[\textbf{c)}] Si $c_n$ son los coeficientes de Fourier de $f$, entonces
    \begin{equation*}
        \sum_{n=-\infty}^{\infty}|c_n|^2\leq \frac{1}{2\pi}\int_{-\pi}^{\pi} f^2(x),
    \end{equation*}
    y, en particular, $c_n\to 0$ cuando $|n|\to\infty$.
\end{enumerate}
\end{tcolorbox}

\textbf{Inciso a).}\,\
Si $n=0$, entonces $e^{inx}=e^0=1$, por lo que
\begin{equation*}
    \frac{1}{2\pi}\int_{-\pi}^{\pi} e^{inx}
    =\frac{1}{2\pi}\int_{-\pi}^{\pi}1
    =\frac{1}{2\pi}\bigl(\pi-(-\pi)\bigr)=1.
\end{equation*}

Si $n\neq0$, entonces una primitiva de $e^{inx}$ es $\frac{e^{inx}}{in}$. Por tanto
\begin{align*}
    \frac{1}{2\pi}\int_{-\pi}^{\pi} e^{inx}
    &=\frac{1}{2\pi}\left[\frac{e^{inx}}{in}\right]_{-\pi}^{\pi} \\
    &=\frac{1}{2\pi in}\left(e^{in\pi}-e^{-in\pi}\right).
\end{align*}
Como $n\in\Z$,
\begin{equation*}
    e^{in\pi}=(-1)^n
    \qquad \text{y}\qquad
    e^{-in\pi}=(-1)^n.
\end{equation*}
Luego
\begin{equation*}
    e^{in\pi}-e^{-in\pi}=0.
\end{equation*}
Así,
\begin{equation*}
    \frac{1}{2\pi}\int_{-\pi}^{\pi} e^{inx}=0.
\end{equation*}
\hfill $\square$

\vspace{0.4cm}

\textbf{Inciso c).}\,\
Definamos el producto interno
\begin{equation*}
    \langle g,h\rangle=\frac{1}{2\pi}\int_{-\pi}^{\pi}g(x)\overline{h(x)}.
\end{equation*}
Por el inciso a), las funciones $e^{inx}$ son ortonormales, es decir,
\begin{equation*}
    \langle e^{inx},e^{imx}\rangle=0\quad(n\neq m),
    \qquad
    \langle e^{inx},e^{inx}\rangle=1.
\end{equation*}
Sea
\begin{equation*}
    S_N(x)=\sum_{n=-N}^{N}c_n e^{inx}.
\end{equation*}
Como $\langle f-S_N,f-S_N\rangle\geq0$, tenemos
\begin{equation*}
    0\leq \frac{1}{2\pi}\int_{-\pi}^{\pi}|f(x)-S_N(x)|^2.
\end{equation*}
Expandiendo el cuadrado usando el producto interno,
\begin{align*}
    0
    &\leq \langle f-S_N,f-S_N\rangle \\
    &=\langle f,f\rangle-\langle f,S_N\rangle-\langle S_N,f\rangle+\langle S_N,S_N\rangle.
\end{align*}
Ahora calculamos cada término. Como
\begin{equation*}
    c_n=\langle f,e^{inx}\rangle,
\end{equation*}
se sigue que
\begin{equation*}
    \langle f,S_N\rangle
    =\left\langle f,\sum_{n=-N}^{N}c_ne^{inx}\right\rangle
    =\sum_{n=-N}^{N}\overline{c_n}\,\langle f,e^{inx}\rangle
    =\sum_{n=-N}^{N}|c_n|^2.
\end{equation*}
Además, por ortonormalidad,
\begin{equation*}
    \langle S_N,S_N\rangle
    =\left\langle \sum_{n=-N}^{N}c_ne^{inx},\sum_{m=-N}^{N}c_me^{imx}\right\rangle
    =\sum_{n=-N}^{N}|c_n|^2.
\end{equation*}
También $\langle S_N,f\rangle=\overline{\langle f,S_N\rangle}=\sum_{n=-N}^{N}|c_n|^2$. Por tanto
\begin{align*}
    0
    &\leq \langle f,f\rangle-\sum_{n=-N}^{N}|c_n|^2-\sum_{n=-N}^{N}|c_n|^2+\sum_{n=-N}^{N}|c_n|^2 \\
    &=\frac{1}{2\pi}\int_{-\pi}^{\pi}f^2(x)-\sum_{n=-N}^{N}|c_n|^2.
\end{align*}
De aquí obtenemos
\begin{equation*}
    \sum_{n=-N}^{N}|c_n|^2
    \leq \frac{1}{2\pi}\int_{-\pi}^{\pi}f^2(x).
\end{equation*}
Como las sumas parciales
\begin{equation*}
    \sum_{n=-N}^{N}|c_n|^2
\end{equation*}
son no decrecientes y están acotadas superiormente, convergen. Por lo tanto
\begin{equation*}
    \sum_{n=-\infty}^{\infty}|c_n|^2
    \leq \frac{1}{2\pi}\int_{-\pi}^{\pi}f^2(x).
\end{equation*}
Finalmente, si una serie de términos no negativos converge, entonces sus términos tienden a $0$. Así
\begin{equation*}
    |c_n|^2\to0
    \qquad\Longrightarrow\qquad
    |c_n|\to0
    \qquad\Longrightarrow\qquad
    c_n\to0
\end{equation*}
cuando $|n|\to\infty$.
\hfill $\square$
\begin{tcolorbox}[title=\textcolor{white}{Problema 4}]
\textit{Sea $f$ una función $2\pi$-periódica con coeficientes complejos de Fourier}
\begin{equation*}
    c_n=\frac{1}{2\pi}\int_{-\pi}^{\pi}f(x)e^{-inx}.
\end{equation*}
\textit{Muestre que si $f$ es par, entonces $c_n=c_{-n}$ y su serie de Fourier sólo contiene términos con cosenos. Muestre también que si $f$ es impar, entonces $c_n=-c_{-n}$ y su serie de Fourier sólo contiene términos con senos.}
\end{tcolorbox}

\textbf{Demostración.}\,\
Primero supongamos que $f$ es par, es decir,
\begin{equation*}
    f(-x)=f(x).
\end{equation*}
Entonces
\begin{align*}
    c_n-c_{-n}
    &=\frac{1}{2\pi}\int_{-\pi}^{\pi}f(x)\bigl(e^{-inx}-e^{inx}\bigr)\\
    &=\frac{1}{2\pi}\int_{-\pi}^{\pi}f(x)(-2i\sen(nx))\\
    &=-\frac{i}{\pi}\int_{-\pi}^{\pi}f(x)\sen(nx).
\end{align*}
Como $f(x)$ es par y $\sen(nx)$ es impar, el producto $f(x)\sen(nx)$ es impar. Por tanto,
\begin{equation*}
    \int_{-\pi}^{\pi}f(x)\sen(nx)=0.
\end{equation*}
Así,
\begin{equation*}
    c_n-c_{-n}=0,
    \qquad \text{es decir,}\qquad c_n=c_{-n}.
\end{equation*}
Por lo tanto, al agrupar los términos $n$ y $-n$ en la serie,
\begin{align*}
    \sum_{n=-\infty}^{\infty}c_ne^{inx}
    &=c_0+\sum_{n=1}^{\infty}\bigl(c_ne^{inx}+c_{-n}e^{-inx}\bigr) \\
    &=c_0+\sum_{n=1}^{\infty}c_n\bigl(e^{inx}+e^{-inx}\bigr) \\
    &=c_0+2\sum_{n=1}^{\infty}c_n\cos(nx).
\end{align*}
Así, en el caso par, sólo aparecen términos con cosenos.

Ahora supongamos que $f$ es impar, es decir,
\begin{equation*}
    f(-x)=-f(x).
\end{equation*}
Entonces
\begin{align*}
    c_n+c_{-n}
    &=\frac{1}{2\pi}\int_{-\pi}^{\pi}f(x)\bigl(e^{-inx}+e^{inx}\bigr)\\
    &=\frac{1}{\pi}\int_{-\pi}^{\pi}f(x)\cos(nx).
\end{align*}
Como $f(x)$ es impar y $\cos(nx)$ es par, el producto $f(x)\cos(nx)$ es impar. Por tanto,
\begin{equation*}
    \int_{-\pi}^{\pi}f(x)\cos(nx)=0.
\end{equation*}
Luego
\begin{equation*}
    c_n+c_{-n}=0,
    \qquad \text{es decir,}\qquad c_{-n}=-c_n.
\end{equation*}
Además $c_0=0$, pues
\begin{equation*}
    c_0=\frac{1}{2\pi}\int_{-\pi}^{\pi}f(x)=0.
\end{equation*}
Entonces, agrupando otra vez los términos $n$ y $-n$,
\begin{align*}
    \sum_{n=-\infty}^{\infty}c_ne^{inx}
    &=\sum_{n=1}^{\infty}\bigl(c_ne^{inx}+c_{-n}e^{-inx}\bigr) \\
    &=\sum_{n=1}^{\infty}c_n\bigl(e^{inx}-e^{-inx}\bigr) \\
    &=2i\sum_{n=1}^{\infty}c_n\sen(nx).
\end{align*}
Así, en el caso impar, sólo aparecen términos con senos.
\hfill $\square$
\begin{tcolorbox}[title=\textcolor{white}{Problema 5}]
\textit{Muestre que, si $f$ es de clase $C^1$ y}
\begin{equation*}
    f(x)\sim \sum_{n=-\infty}^{\infty}c_ne^{inx},
\end{equation*}
\textit{entonces}
\begin{equation*}
    f'(x)\sim \sum_{n=-\infty}^{\infty}inc_ne^{inx}.
\end{equation*}
\end{tcolorbox}

\textbf{Demostración.}\,\
Sea $d_n$ el $n$-ésimo coeficiente de Fourier de $f'$. Entonces
\begin{equation*}
    d_n=\frac{1}{2\pi}\int_{-\pi}^{\pi}f'(x)e^{-inx}.
\end{equation*}
Integramos por partes. Tomamos
\begin{equation*}
    u=e^{-inx},\qquad dv=f'(x)\,dx.
\end{equation*}
Entonces
\begin{equation*}
    du=-in e^{-inx}\,dx,
    \qquad
    v=f(x).
\end{equation*}
Por tanto
\begin{align*}
    d_n
    &=\frac{1}{2\pi}\left(\left[f(x)e^{-inx}\right]_{-\pi}^{\pi}
    +in\int_{-\pi}^{\pi}f(x)e^{-inx}\right).
\end{align*}
Como $f$ es $2\pi$-periódica, se tiene $f(\pi)=f(-\pi)$. Además, como $n\in\Z$,
\begin{equation*}
    e^{-in\pi}=e^{in\pi}=(-1)^n.
\end{equation*}
Por eso el término de frontera se anula:
\begin{align*}
    \left[f(x)e^{-inx}\right]_{-\pi}^{\pi}
    &=f(\pi)e^{-in\pi}-f(-\pi)e^{in\pi} \\
    &=f(\pi)(-1)^n-f(-\pi)(-1)^n \\
    &=0.
\end{align*}
Luego
\begin{align*}
    d_n
    &=\frac{in}{2\pi}\int_{-\pi}^{\pi}f(x)e^{-inx} \\
    &=in c_n.
\end{align*}
Esto significa que los coeficientes de Fourier de $f'$ son precisamente $inc_n$. Por tanto
\begin{equation*}
    f'(x)\sim \sum_{n=-\infty}^{\infty}inc_ne^{inx}.
\end{equation*}
\hfill $\square$

\vspace{0.4cm}
\begin{tcolorbox}[title=\textcolor{white}{Problema 6(a)-(b)}]
\textit{Considere la función}
\begin{equation*}
    f(x)=|x|,
\end{equation*}
\textit{definida en el intervalo $[-\pi,\pi]$, y extiéndala a toda la recta real de tal manera que sea periódica de periodo $2\pi$.}
\begin{enumerate}
    \item[\textbf{a)}] \textit{Calcule los coeficientes de Fourier de $f$ usando integración por partes.}
    \item[\textbf{b)}] \textit{Escriba la serie de Fourier de $f$ en términos de senos y cosenos.}
\end{enumerate}
\end{tcolorbox}

\textbf{Demostración.}\,\\
Recordemos que
\begin{equation*}
    c_n=\frac{1}{2\pi}\int_{-\pi}^{\pi}|x|e^{-inx}\dd.
\end{equation*}
Como $|x|$ es una función par, al escribir
\begin{equation*}
    e^{-inx}=\cos(nx)-i\sen(nx)
\end{equation*}
observamos que $|x|\sen(nx)$ es impar. Por tanto,
\begin{align*}
    c_n
    &=\frac{1}{2\pi}\int_{-\pi}^{\pi}|x|\bigl(\cos(nx)-i\sen(nx)\bigr)\dd \\
    &=\frac{1}{2\pi}\int_{-\pi}^{\pi}|x|\cos(nx)\dd \\
    &=\frac{1}{\pi}\int_{0}^{\pi}x\cos(nx)\dd,
\end{align*}
para $n\neq 0$.

Integramos por partes tomando
\begin{equation*}
    u=x,\qquad dv=\cos(nx)\,dx,
\end{equation*}
de modo que
\begin{equation*}
    du=dx,\qquad v=\frac{\sen(nx)}{n}.
\end{equation*}
Entonces
\begin{align*}
    \int_{0}^{\pi}x\cos(nx)\dd
    &=\left.\frac{x\sen(nx)}{n}\right|{0}^{\pi}-\frac{1}{n}\int{0}^{\pi}\sen(nx)\dd \\
    &=0-\frac{1}{n}\left[-\frac{\cos(nx)}{n}\right]_{0}^{\pi} \\
    &=\frac{1}{n^2}\bigl(\cos(n\pi)-1\bigr) \\
    &=\frac{(-1)^n-1}{n^2}.
\end{align*}
Por consiguiente, para $n\neq 0$ se tiene
\begin{equation*}
    c_n=\frac{(-1)^n-1}{\pi n^2}.
\end{equation*}

Para $n=0$,
\begin{align*}
    c_0
    &=\frac{1}{2\pi}\int_{-\pi}^{\pi}|x|\dd
    =\frac{1}{\pi}\int_{0}^{\pi}x\dd
    =\frac{1}{\pi}\left.\frac{x^2}{2}\right|_{0}^{\pi}
    =\frac{\pi}{2}.
\end{align*}

Así, la serie de Fourier compleja de $|x|$ es
\begin{equation*}
    |x|\sim \sum_{n=-\infty}^{\infty}c_ne^{inx},
\end{equation*}
con
\begin{equation*}
    c_0=\frac{\pi}{2},
    \qquad
    c_n=\frac{(-1)^n-1}{\pi n^2}\quad (n\neq 0).
\end{equation*}
Es decir,
\begin{equation*}
    |x|\sim \frac{\pi}{2}+\sum_{n\neq 0}\frac{(-1)^n-1}{\pi n^2}e^{inx}.
\end{equation*}
Como los coeficientes son reales y satisfacen $c_{-n}=c_n$, esta serie también puede escribirse como
\begin{equation*}
    |x|\sim \frac{\pi}{2}+\frac{2}{\pi}\sum_{n=1}^{\infty}\frac{(-1)^n-1}{n^2}\cos(nx)
\end{equation*}
\hfill $\square$
\vspace{0.4cm}
\begin{tcolorbox}[title=\textcolor{white}{Problema 7}]
\textit{Pruebe que, si $f$ es de clase $C^1$, entonces para cada punto $x$ interior del dominio existen $\delta>0$ y $M>0$ tales que}
\begin{equation*}
    |f(x+t)-f(x)|\leq M|t|
    \qquad\text{para todo } |t|<\delta.
\end{equation*}
\end{tcolorbox}

\textbf{Demostración.}\,\
Fijemos un punto $x$ interior del dominio de $f$. Como $x$ es interior, existe $\delta_0>0$ tal que
\begin{equation*}
    [x-\delta_0,x+\delta_0]
\end{equation*}
está contenido en el dominio de $f$.

Como $f$ es de clase $C^1$, su derivada $f'$ es continua. En particular, $f'$ es continua en el intervalo compacto
\begin{equation*}
    [x-\delta_0,x+\delta_0].
\end{equation*}
Por el teorema del valor extremo, $|f'|$ alcanza un máximo en ese intervalo. Es decir, existe $M>0$ tal que
\begin{equation*}
    |f'(y)|\leq M
    \qquad\text{para todo }y\in[x-\delta_0,x+\delta_0].
\end{equation*}
Tomemos $\delta=\delta_0$. Si $|t|<\delta$, entonces $x+t\in[x-\delta,x+\delta]$.

Aplicando el Teorema del Valor Medio a $f$ en el intervalo cuyos extremos son $x$ y $x+t$, existe un punto $c$ entre $x$ y $x+t$ tal que
\begin{equation*}
    f(x+t)-f(x)=f'(c)t.
\end{equation*}
Por lo tanto,
\begin{equation*}
    |f(x+t)-f(x)|=|f'(c)|\,|t|\leq M|t|.
\end{equation*}
Así existen $\delta>0$ y $M>0$ que satisfacen la desigualdad pedida.
Ahora probaremos que $\displaystyle \sum_{n\in \N}\frac{1}{n^2}=\frac{\pi}{6}$, promero supongamos 
primero supondremos que  la serie de $f(x)=|x|$ converge puntualmente en $x=0$, luego tenemos que:\,\\
\begin{equation*}
    \frac{\pi}{2}+\frac{1}{\pi}\sum_{-\infty, n\,\,impar}^{\infty}\frac{(-1)^n-1}{n^2}=0
\end{equation*} 
\begin{equation*}
    \frac{\pi}{2}-\frac{4}{\pi}\sum_{n\,\,\text{impar}}\frac{1}{n^2}=0\implies \sum_{n\text{impar}}\frac{1}{n^2}=\frac{\pi^2}{8}
\end{equation*}\,\\
luego tenemos que:\,\\
\begin{equation*}
    \sum_{n\in \N}\frac{1}{n^2}=\sum_{n\,\text{impar}}\frac{1}{n^2}+\sum_{n\,\text{par}}\frac{1}{n^2}
\end{equation*}\,\\
\begin{equation*}
     \sum_{n\in \N}\frac{1}{n^2}=\sum_{n\,\text{impar}}\frac{1}{n^2}+\frac{1}{4}\sum_{n\,\text{par}}\frac{1}{n^2}
\end{equation*}\,\\
\begin{equation*}
    \frac{3}{4}\sum_{n\in N}\frac{1}{n^2}=\frac{\pi^2}{8}\implies \sum_{n\in\N}\frac{1}{n^2}=\frac{\pi^2}{6}
\end{equation*}


\hfill $\square$
\begin{tcolorbox}[title=\textcolor{white}{Problema 8}]
\textit{Sea}
\begin{equation*}
    D_N(x)=\sum_{n=-N}^{N}e^{inx}.
\end{equation*}
\textit{Pruebe las siguientes propiedades del núcleo de Dirichlet:}
\begin{enumerate}
    \item[\textbf{a)}] Para $x\notin 2\pi\Z$,
    \begin{equation*}
        D_N(x)=\frac{\sen\left(\left(N+\frac12\right)x\right)}{\sen\left(\frac{x}{2}\right)}.
    \end{equation*}
    \item[\textbf{b)}] $\displaystyle \frac{1}{2\pi}\int_{-\pi}^{\pi}D_N(t)dt=1$.
    \item[\textbf{c)}] Si $S_N(f,x)=\displaystyle\sum_{n=-N}^{N}c_ne^{inx}$, entonces
    \begin{equation*}
        \frac{1}{2\pi}\int_{-\pi}^{\pi}f(t)D_N(x-t)dt=S_N(f,x).
    \end{equation*}
    \item[\textbf{d)}]
    \begin{equation*}
        \frac{1}{2\pi}\int_{-\pi}^{\pi}f(t)D_N(x-t)dt
        =\frac{1}{2\pi}\int_{-\pi}^{\pi}f(x-t)D_N(t)dt.
    \end{equation*}
\end{enumerate}
\end{tcolorbox}

\textbf{Inciso a).}\,\
Partimos de
\begin{equation*}
    D_N(x)=\sum_{n=-N}^{N}e^{inx}.
\end{equation*}
Multiplicamos por $e^{ix}-1$:
\begin{align*}
    (e^{ix}-1)D_N(x)
    &=(e^{ix}-1)\sum_{n=-N}^{N}e^{inx} \\
    &=\sum_{n=-N}^{N}\bigl(e^{i(n+1)x}-e^{inx}\bigr).
\end{align*}
La suma es telescópica, así que
\begin{equation*}
    (e^{ix}-1)D_N(x)=e^{i(N+1)x}-e^{-iNx}.
\end{equation*}
Multiplicamos numerador y denominador por $e^{-ix/2}$:
\begin{align*}
    D_N(x)
    &=\frac{e^{i(N+1)x}-e^{-iNx}}{e^{ix}-1} \\
    &=\frac{e^{i\left(N+\frac12\right)x}-e^{-i\left(N+\frac12\right)x}}{e^{ix/2}-e^{-ix/2}}.
\end{align*}
Usando
\begin{equation*}
    e^{ia}-e^{-ia}=2i\sen(a),
\end{equation*}
obtenemos
\begin{equation*}
    D_N(x)=\frac{\sen\left(\left(N+\frac12\right)x\right)}{\sen\left(\frac{x}{2}\right)}.
\end{equation*}
Además, directamente de la definición,
\begin{equation*}
    D_N(0)=\sum_{n=-N}^{N}1=2N+1.
\end{equation*}

\vspace{0.3cm}

\textbf{Inciso b).}\,\
Por linealidad de la integral,
\begin{align*}
    \frac{1}{2\pi}\int_{-\pi}^{\pi}D_N(t)dt
    &=\frac{1}{2\pi}\int_{-\pi}^{\pi}\sum_{n=-N}^{N}e^{int}dt \\
    &=\sum_{n=-N}^{N}\frac{1}{2\pi}\int_{-\pi}^{\pi}e^{int}dt.
\end{align*}
Por el Problema 3, todos los términos con $n\neq0$ valen $0$, mientras que el término con $n=0$ vale $1$. Por tanto,
\begin{equation*}
    \frac{1}{2\pi}\int_{-\pi}^{\pi}D_N(t)dt=1.
\end{equation*}

\vspace{0.3cm}

\textbf{Inciso c).}\,\
Usando la definición de $D_N$,
\begin{align*}
    \frac{1}{2\pi}\int_{-\pi}^{\pi}f(t)D_N(x-t)dt
    &=\frac{1}{2\pi}\int_{-\pi}^{\pi}f(t)\sum_{n=-N}^{N}e^{in(x-t)}dt \\
    &=\sum_{n=-N}^{N}e^{inx}\frac{1}{2\pi}\int_{-\pi}^{\pi}f(t)e^{-int}dt \\
    &=\sum_{n=-N}^{N}c_ne^{inx} \\
    &=S_N(f,x).
\end{align*}

\vspace{0.3cm}

\textbf{Inciso d).}\,\
Hacemos el cambio de variable
\begin{equation*}
    u=x-t,
    \qquad du=-dt.
\end{equation*}
Entonces
\begin{equation*}
    \frac{1}{2\pi}\int_{-\pi}^{\pi}f(t)D_N(x-t)dt
    =\frac{1}{2\pi}\int_{x-\pi}^{x+\pi}f(x-u)D_N(u)du.
\end{equation*}
Como $f$ y $D_N$ son $2\pi$-periódicas, el producto $f(x-u)D_N(u)$ también es $2\pi$-periódico. Por el Problema 2, podemos integrar en cualquier intervalo de longitud $2\pi$. Así,
\begin{equation*}
    \frac{1}{2\pi}\int_{x-\pi}^{x+\pi}f(x-u)D_N(u)du
    =\frac{1}{2\pi}\int_{-\pi}^{\pi}f(x-u)D_N(u)du.
\end{equation*}
Renombrando $u$ como $t$, queda
\begin{equation*}
    \frac{1}{2\pi}\int_{-\pi}^{\pi}f(t)D_N(x-t)dt 
    =\frac{1}{2\pi}\int_{-\pi}^{\pi}f(x-t)D_N(t)dt.
\end{equation*}
\hfill $\square$
\begin{tcolorbox}[title=\textcolor{white}{Problema 9}]
\textit{Sea}
\begin{equation*}
    g(t)=
    \begin{cases}
    \displaystyle \frac{f(x-t)-f(x)}{\sen\left(\frac{t}{2}\right)}, & 0<|t|\leq \pi,\\[0.3cm]
    0, & t=0.
    \end{cases}
\end{equation*}
\textit{Muestre que}
\begin{align*}
    S_N(f,x)-f(x)
    &=\frac{1}{2\pi}\int_{-\pi}^{\pi}g(t)\sen\left(\left(N+\frac12\right)t\right)dt \\
    &=\frac{1}{2\pi}\int_{-\pi}^{\pi}g(t)\cos\left(\frac{t}{2}\right)\sen(Nt)dt
    +\frac{1}{2\pi}\int_{-\pi}^{\pi}g(t)\sen\left(\frac{t}{2}\right)\cos(Nt)dt.
\end{align*}
\end{tcolorbox}

\textbf{Demostración.}\,\
Por el Problema 8, incisos b) y d), tenemos
\begin{align*}
    S_N(f,x)-f(x)
    &=\frac{1}{2\pi}\int_{-\pi}^{\pi}f(x-t)D_N(t)dt
      -f(x)\frac{1}{2\pi}\int_{-\pi}^{\pi}D_N(t)dt \\
    &=\frac{1}{2\pi}\int_{-\pi}^{\pi}\bigl(f(x-t)-f(x)\bigr)D_N(t)dt.
\end{align*}
Ahora usamos la fórmula del núcleo de Dirichlet:
\begin{equation*}
    D_N(t)=\frac{\sen\left(\left(N+\frac12\right)t\right)}{\sen\left(\frac{t}{2}\right)}.
\end{equation*}
Entonces
\begin{align*}
    S_N(f,x)-f(x)
    &=\frac{1}{2\pi}\int_{-\pi}^{\pi}
    \frac{f(x-t)-f(x)}{\sen\left(\frac{t}{2}\right)}
    \sen\left(\left(N+\frac12\right)t\right)dt \\
    &=\frac{1}{2\pi}\int_{-\pi}^{\pi}g(t)\sen\left(\left(N+\frac12\right)t\right)dt.
\end{align*}
El valor de $g$ en $t=0$ no afecta la integral.

Finalmente usamos la identidad trigonométrica
\begin{equation*}
    \sen(\alpha+\beta)=\sen(\alpha)\cos(\beta)+\sen(\beta)\cos(\alpha).
\end{equation*}
Tomando $\alpha=Nt$ y $\beta=t/2$, obtenemos
\begin{equation*}
    \sen\left(\left(N+\frac12\right)t\right)
    =\sen(Nt)\cos\left(\frac{t}{2}\right)
    +\sen\left(\frac{t}{2}\right)\cos(Nt).
\end{equation*}
Sustituyendo en la integral,
\begin{align*}
    S_N(f,x)-f(x)
    &=\frac{1}{2\pi}\int_{-\pi}^{\pi}g(t)\cos\left(\frac{t}{2}\right)\sen(Nt)dt \\
    &\quad +\frac{1}{2\pi}\int_{-\pi}^{\pi}g(t)\sen\left(\frac{t}{2}\right)\cos(Nt)dt.
\end{align*}
\hfill $\square$
\begin{tcolorbox}[title=\textcolor{white}{Problema 10}]
\textit{Prueba que $g(t)cos\left(\frac{t}{2}\right)$ y $g(t)sen\left(\frac{t}{2}\right)$ est\'an acotadas en $[-\pi,\pi]$ y que son Rieman integrables}
\end{tcolorbox}
\textbf{Demostraci\'on}: Primero veamos que $f$ es continua sea $\epsilon_0>0$ tomemos $\delta_1:=min\left\{\delta,\frac{\epsilon_0}{M}\right\}$, entonces si $|x|\leq \delta$:\,\\
\begin{equation*}
    |f(x+t)-f(x)|\leq M|t|\leq \frac{M \epsilon}{M}=\epsilon
\end{equation*}
Por tanto $\displaystyle \lim_{t\to 0}f(x+t)=f(x)$, ie $f$ es continua, ahora sea
\[
f(x) =
\begin{cases}
\frac{f(x-t)-f(x)}{sen\left(\frac{t}{2}\right)} \quad & \text{si } 0<|t|\leq\pi, \\
0 \quad  & \text{si } x =0.
\end{cases}
\]
Demostraremos que $g$ es acotada, veamos que:\,\\
\begin{equation*}
    \lim_{t\to 1} \frac{sen\left(\frac{t}{2}\right)}{t}=\frac{1}{2}
\end{equation*}
por tanto dado $\epsilon=\frac{1}{4}>0$ existe $\delta_0>0$ tal que si $0<|t|\leq \delta$
\begin{equation*}
    \left|\frac{sen\left(\frac{t}{2}\right)}{t}-\frac{1}{2}\right|\leq \frac{1}{4}
\end{equation*}\,\\
\begin{equation*}
    -\frac{1}{4}\leq\frac{sen\left(\frac{t}{2}\right)}{t}-\frac{1}{2}\leq \frac{1}{4}\implies
\end{equation*}
\begin{equation*}
    \frac{1}{4}\leq\frac{sen\left(\frac{t}{2}\right)}{t}\leq \frac{3}{4}\implies
\end{equation*}
\begin{equation*}
    \frac{4}{3}\leq\frac{|t|}{\left|sen\left(\frac{t}{2}\right)\right|}\leq 4
\end{equation*}
Tomemos $\delta_2=min\{\delta,\delta_0\}$, entonces si $0<|t|\leq\delta_2$:\,\\
\begin{equation*}
    \left|\frac{f(x-t)-f(x)}{sen\left(\frac{t}{2}\right)}\right|\leq \frac{|t|M}{\left|sen\left(\frac{t}{2}\right)\right|}
\end{equation*}\,\\
Luego $g$ esta acotada para cada $x\in (-\delta_2,\delta_2)$, ahora en los intervalos $[-\pi,-\delta_2]$ y $[\delta_2,\pi]$ la funci\'on es continua y por tanto acotada,
si tomamos el m\'aximo de estas tres cotas tenemos que $g$ es acotada en su dominio, por tanto $g(t)cos(t)$ y $g(t)sen(t)$ tambien lo son, ahora para ver que son integrables basta notar que $g$ solo puede ser discontinua en $0$, como su conjunto de discontinuidades es finito en $[-\pi,\pi]$ por teorema de Lebesgue se sigue que $g$ es integrable y por tanto $g(t)cos(t/2),g(t)sen(t/2)$ son integrables 
\hfill $\square$

\begin{tcolorbox}[title=\textcolor{white}{Problema 11}]
\textit{Deduzca que ambas integrales que aparecen en la identidad del Problema 9 van a cero cuando $N\to\infty$. Use que los coeficientes de Fourier de cualquier función Riemann integrable tienden a cero, como se probó en el Problema 3(c). Esto completa la prueba del Teorema 1.}
\end{tcolorbox}

\textbf{Demostración.}\,\
Por el Problema 9,
\begin{align*}
    S_N(f,x)-f(x)
    &=\frac{1}{2\pi}\int_{-\pi}^{\pi}g(t)\cos\left(\frac{t}{2}\right)\sen(Nt)\dt \\
    &\quad +\frac{1}{2\pi}\int_{-\pi}^{\pi}g(t)\sen\left(\frac{t}{2}\right)\cos(Nt)\dt.
\end{align*}
Definimos
\begin{equation*}
    h_1(t)=g(t)\cos\left(\frac{t}{2}\right),
    \qquad
    h_2(t)=g(t)\sen\left(\frac{t}{2}\right).
\end{equation*}
Como $g$ es Riemann integrable y $\cos(t/2)$, $\sen(t/2)$ son continuas, entonces $h_1$ y $h_2$ son Riemann integrables.

Recordemos que si $h$ es Riemann integrable y
\begin{equation*}
    \widehat{h}n=\frac{1}{2\pi}\int{-\pi}^{\pi}h(t)e^{-int}\dt,
\end{equation*}
entonces, por el resultado de los coeficientes de Fourier,
\begin{equation*}
    \widehat{h}_n\to0.
\end{equation*}
Pero
\begin{equation*}
    \widehat{h}_n
    =\frac{1}{2\pi}\int_{-\pi}^{\pi}h(t)\cos(nt)\dt
    -i\frac{1}{2\pi}\int_{-\pi}^{\pi}h(t)\sen(nt)\dt.
\end{equation*}
Por tanto, tanto la parte real como la parte imaginaria tienden a $0$. Es decir,
\begin{equation*}
    \frac{1}{2\pi}\int_{-\pi}^{\pi}h(t)\cos(nt)\dt\to0,
    \qquad
    \frac{1}{2\pi}\int_{-\pi}^{\pi}h(t)\sen(nt)\dt\to0.
\end{equation*}
Aplicando esto a $h_1$ y $h_2$, obtenemos
\begin{equation*}
    \frac{1}{2\pi}\int_{-\pi}^{\pi}h_1(t)\sen(Nt)\dt\to0,
    \qquad
    \frac{1}{2\pi}\int_{-\pi}^{\pi}h_2(t)\cos(Nt)\dt\to0.
\end{equation*}
Así,
\begin{equation*}
    S_N(f,x)-f(x)\to0.
\end{equation*}
Por lo tanto,
\begin{equation*}
    \lim_{N\to\infty}S_N(f,x)=f(x).
\end{equation*}
\hfill $\square$

\newpage

\begin{tcolorbox}[title=\textcolor{white}{Problema 12}]
\textit{Suponga que $\gamma:\R\to\mathbb C$, con}
\begin{equation*}
    \gamma(t)=x(t)+iy(t),
\end{equation*}
\textit{es una curva cerrada de longitud $L$ y rapidez constante}
\begin{equation*}
    |\gamma'(t)|=\frac{L}{2\pi}.
\end{equation*}
\textit{Sean $c_n$, con $n\in\Z$, los coeficientes de Fourier de $\gamma$:}
\begin{equation*}
    c_n=\frac{1}{2\pi}\int_{-\pi}^{\pi}\gamma(t)e^{-int}\,dt.
\end{equation*}
\textit{Muestre que}
\begin{equation*}
    \gamma(t)=\sum_{n\in\Z}c_ne^{int}
    \qquad\text{y}\qquad
    \gamma'(t)=i\sum_{n\in\Z}nc_ne^{int}.
\end{equation*}
\end{tcolorbox}

\textbf{Demostración.}\,\
Escribimos la curva en sus partes real e imaginaria:
\begin{equation*}
    \gamma(t)=x(t)+iy(t).
\end{equation*}
Entonces
\begin{align*}
    C_n
    &=\frac{1}{2\pi}\int_{-\pi}^{\pi}\gamma(t)e^{-int}\dt \\
    &=\frac{1}{2\pi}\int_{-\pi}^{\pi}\bigl(x(t)+iy(t)\bigr)e^{-int}\dt \\
    &=\frac{1}{2\pi}\int_{-\pi}^{\pi}x(t)e^{-int}\dt
    +i\frac{1}{2\pi}\int_{-\pi}^{\pi}y(t)e^{-int}\dt.
\end{align*}
Definimos
\begin{equation*}
    C_n^x=\frac{1}{2\pi}\int_{-\pi}^{\pi}x(t)e^{-int}\dt,
    \qquad
    C_n^y=\frac{1}{2\pi}\int_{-\pi}^{\pi}y(t)e^{-int}\dt.
\end{equation*}
Así,
\begin{equation*}
    C_n=C_n^x+iC_n^y.
\end{equation*}
Ahora consideremos la suma parcial de Fourier de $\gamma$:
\begin{align*}
    \sum_{n=-N}^{N}C_ne^{int}
    &=\sum_{n=-N}^{N}\bigl(C_n^x+iC_n^y\bigr)e^{int} \\
    &=\sum_{n=-N}^{N}C_n^x e^{int}
    +i\sum_{n=-N}^{N}C_n^y e^{int} \\
    &=S_N(x,t)+iS_N(y,t).
\end{align*}
Como la curva es suave, sus coordenadas $x$ y $y$ cumplen las hipótesis del teorema de convergencia usado para series de Fourier. Por tanto,
\begin{equation*}
    \lim_{N\to\infty}S_N(x,t)=x(t),
    \qquad
    \lim_{N\to\infty}S_N(y,t)=y(t).
\end{equation*}
Así,\begin{align*}
    \lim_{N\to\infty}\sum_{n=-N}^{N}C_ne^{int}
    &=x(t)+iy(t) \\
    &=\gamma(t).
\end{align*}
Por lo tanto,
\begin{equation*}
    \gamma(t)\sim\sum_{n\in\Z}C_ne^{int}.
\end{equation*}
\hfill $\square$

\newpage

\begin{tcolorbox}[title=\textcolor{white}{Problema 13}]
\textit{Ahora necesitamos un teorema de cálculo integral: el área de la región encerrada por la curva $\gamma(t)=(x(t),y(t))$ está dada por}
\begin{equation*}
    A=\frac12\int_{-\pi}^{\pi}\bigl(x(t)y'(t)-x'(t)y(t)\bigr)\,dt.
\end{equation*}
\textit{Suponga, además, que la curva se recorre de tal manera que la región encerrada siempre queda a su izquierda. Resuelva esta integral para las curvas}
\begin{equation*}
    \alpha(t)=e^{it}
    \qquad\text{y}\qquad
    \beta(t)=e^{-it}.
\end{equation*}
\textit{Al final del problema se dará una demostración de esta fórmula.}
\end{tcolorbox}

\textbf{Demostración.}\,
Usamos la fórmula del área orientada
\begin{equation*}
    A=\frac12\int_{-\pi}^{\pi}\bigl(x(t)y'(t)-x'(t)y(t)\bigr)\dt.
\end{equation*}

Para $\gamma_1(t)=\cos t+i\sen t$ tenemos
\begin{equation*}
    x_1(t)=\cos t,
    \qquad
    y_1(t)=\sen t,
    \qquad
    x_1'(t)=-\sen t,
    \qquad
    y_1'(t)=\cos t.
\end{equation*}
Entonces
\begin{align*}
    A_1
    &=\frac12\int_{-\pi}^{\pi}\bigl(x_1(t)y_1'(t)-x_1'(t)y_1(t)\bigr)\dt \\
    &=\frac12\int_{-\pi}^{\pi}\bigl(\cos^2 t+\sen^2 t\bigr)\dt \\
    &=\frac12\int_{-\pi}^{\pi}1\,\dt \\
    &=\frac12(2\pi)=\pi.
\end{align*}

Ahora, para $\gamma_2(t)=\cos t-i\sen t$ tenemos
\begin{equation*}
    x_2(t)=\cos t,
    \qquad
    y_2(t)=-\sen t,
    \qquad
    x_2'(t)=-\sen t,
    \qquad
    y_2'(t)=-\cos t.
\end{equation*}
Por tanto,
\begin{align*}
    A_2
    &=\frac12\int_{-\pi}^{\pi}\bigl(x_2(t)y_2'(t)-x_2'(t)y_2(t)\bigr)\dt \\
    &=\frac12\int_{-\pi}^{\pi}\bigl(-\cos^2 t-\sen^2 t\bigr)\dt \\
    &=\frac12\int_{-\pi}^{\pi}(-1)\,\dt \\
    &=-\frac12(2\pi)=-\pi.
\end{align*}
Así, las áreas orientadas son
\begin{equation*}
    A_1=\pi,
    \qquad
    A_2=-\pi.\end{equation*}
\hfill $\square$

\newpage

\begin{tcolorbox}[title=\textcolor{white}{Problema 14}]
\textit{Observe que}
\begin{equation*}
    x(t)y'(t)-x'(t)y(t)=\Im\bigl(\overline{\gamma(t)}\,\gamma'(t)\bigr),
\end{equation*}
\textit{donde $\Im(z)$ denota la parte imaginaria del número complejo $z$. Pruebe que}
\begin{equation*}
    A=\pi\sum_{n\in\Z}n|c_n|^2.
\end{equation*}
\end{tcolorbox}

\textbf{Demostración.}\,
Recordemos que el área orientada puede escribirse como
\begin{equation*}
    A=\frac12\int_{-\pi}^{\pi}\Im\bigl(\overline{\gamma(t)}\,\gamma'(t)\bigr)\dt.
\end{equation*}
Además,
\begin{equation*}
    \gamma'(t)=i\sum_{n\in\Z}nc_ne^{int},
    \qquad
    \overline{\gamma(t)}=\sum_{k\in\Z}\overline{c_k}e^{-ikt}.
\end{equation*}
Sustituyendo,
\begin{align*}
    \overline{\gamma(t)}\,\gamma'(t)
    &=\left(\sum_{k\in\Z}\overline{c_k}e^{-ikt}\right)
      \left(i\sum_{n\in\Z}nc_ne^{int}\right) \\
    &=i\sum_{n\in\Z}\sum_{k\in\Z}n c_n\overline{c_k}\,e^{i(n-k)t}.
\end{align*}
Integrando en $[-\pi,\pi]$ y usando la ortogonalidad de las exponenciales complejas,
\begin{equation*}
    \int_{-\pi}^{\pi}e^{i(n-k)t}\dt=
    \begin{cases}
        2\pi, & n=k,\\
        0, & n\neq k,
    \end{cases}
\end{equation*}
obtenemos
\begin{align*}
    \int_{-\pi}^{\pi}\overline{\gamma(t)}\,\gamma'(t)\dt
    &=i\sum_{n,k}n c_n\overline{c_k}\int_{-\pi}^{\pi}e^{i(n-k)t}\dt \\
    &=2\pi i\sum_{n\in\Z}n c_n\overline{c_n} \\
    &=2\pi i\sum_{n\in\Z}n|c_n|^2.
\end{align*}
El término de la derecha es puramente imaginario, por lo que
\begin{equation*}
    \Im\left(\int_{-\pi}^{\pi}\overline{\gamma(t)}\,\gamma'(t)\dt\right)
    =2\pi\sum_{n\in\Z}n|c_n|^2.
\end{equation*}
Finalmente,
\begin{equation*}
    A=\frac12\cdot 2\pi\sum_{n\in\Z}n|c_n|^2
    =\pi\sum_{n\in\Z}n|c_n|^2.
\end{equation*}
\hfill $\square$

\newpage

\begin{tcolorbox}[title=\textcolor{white}{Problema 15}]
\textit{Pruebe que}
\begin{equation*}
    L^2=4\pi^2\sum_{n\in\Z}n^2|c_n|^2.
\end{equation*}
\end{tcolorbox}

\textbf{Demostración.}\,
Como la curva está parametrizada con rapidez constante,
\begin{equation*}
    |\gamma'(t)|=\frac{L}{2\pi}
    \qquad\text{para todo } t\in[-\pi,\pi].
\end{equation*}
Entonces
\begin{align*}
    \int_{-\pi}^{\pi}|\gamma'(t)|^2\dt
    &=\int_{-\pi}^{\pi}\left(\frac{L}{2\pi}\right)^2\dt \\
    &=2\pi\cdot\frac{L^2}{4\pi^2} \\
    &=\frac{L^2}{2\pi}.
\end{align*}

Por otra parte, al derivar la serie de Fourier de $\gamma$ obtenemos
\begin{equation*}
    \gamma'(t)=i\sum_{n\in\Z}n c_n e^{int}.
\end{equation*}
Aplicando la identidad de Parseval a $\gamma'$, cuyos coeficientes de Fourier son $inc_n$, resulta
\begin{equation*}
    \frac{1}{2\pi}\int_{-\pi}^{\pi}|\gamma'(t)|^2\dt
    =\sum_{n\in\Z}|inc_n|^2
    =\sum_{n\in\Z}n^2|c_n|^2.
\end{equation*}
Multiplicando por $2\pi$,
\begin{equation*}
    \int_{-\pi}^{\pi}|\gamma'(t)|^2\dt
    =2\pi\sum_{n\in\Z}n^2|c_n|^2.
\end{equation*}
Comparando con el valor calculado antes,
\begin{equation*}
    \frac{L^2}{2\pi}=2\pi\sum_{n\in\Z}n^2|c_n|^2.
\end{equation*}
Por lo tanto,\begin{equation*}
    L^2=4\pi^2\sum_{n\in\Z}n^2|c_n|^2.
\end{equation*}
\hfill $\square$

\newpage

\begin{tcolorbox}[title=\textcolor{white}{Problema 16}]
\textit{Concluya que}
\begin{equation*}
    L^2-4\pi A\geq 0,
\end{equation*}
\textit{con igualdad si y sólo si}
\begin{equation*}
    \gamma(t)=c_0+c_1e^{it},
    \qquad |c_1|=\frac{L}{2\pi}.
\end{equation*}
\textit{En este caso, la curva es un círculo con centro en $c_0$ y radio $|c_1|$.}
\end{tcolorbox}

\textbf{Demostración.}\,
Por los problemas anteriores,
\begin{equation*}
    L^2=4\pi^2\sum_{n\in\Z}n^2|c_n|^2,
    \qquad
    A=\pi\sum_{n\in\Z}n|c_n|^2.
\end{equation*}
Por tanto,
\begin{align*}
    L^2-4\pi A
    &=4\pi^2\sum_{n\in\Z}n^2|c_n|^2
      -4\pi^2\sum_{n\in\Z}n|c_n|^2 \\
    &=4\pi^2\sum_{n\in\Z}\bigl(n^2-n\bigr)|c_n|^2 \\
    &=4\pi^2\sum_{n\in\Z}n(n-1)|c_n|^2.
\end{align*}
Como para todo entero $n$ se cumple $n(n-1)\ge 0$, deducimos que cada sumando es no negativo y, en consecuencia,
\begin{equation*}
    L^2-4\pi A\ge 0.
\end{equation*}

Ahora estudiemos el caso de igualdad. Se tiene
\begin{equation*}
    L^2-4\pi A=0
    \iff
    \sum_{n\in\Z}n(n-1)|c_n|^2=0.
\end{equation*}
Como todos los términos de la suma son no negativos, esto sólo puede ocurrir si
\begin{equation*}
    n(n-1)|c_n|^2=0
    \qquad\text{para todo } n\in\Z.
\end{equation*}
Así, necesariamente
\begin{equation*}
    c_n=0
    \qquad\text{para todo } n\neq 0,1.
\end{equation*}
Por lo tanto,
\begin{equation*}
    \gamma(t)=c_0+c_1e^{it}.
\end{equation*}
Además,
\begin{equation*}
    \gamma'(t)=ic_1e^{it},
\qquad\text{de donde}\qquad
    |\gamma'(t)|=|c_1|.
\end{equation*}
Como la rapidez constante es $L/(2\pi)$, se sigue que
\begin{equation*}
    |c_1|=\frac{L}{2\pi}.
\end{equation*}

Recíprocamente, si
\begin{equation*}
    \gamma(t)=c_0+c_1e^{it}
    \qquad\text{con}\qquad
    |c_1|=\frac{L}{2\pi},
\end{equation*}
entonces $c_n=0$ para todo $n\neq 0,1$, y por ello
\begin{align*}
    L^2-4\pi A
    &=4\pi^2\sum_{n\in\Z}n(n-1)|c_n|^2 \\
    &=4\pi^2\Bigl(0\cdot(-1)|c_0|^2+1\cdot 0\,|c_1|^2\Bigr) \\
    &=0.
\end{align*}
Con esto queda probado que
\begin{equation*}
    L^2-4\pi A\ge 0,
\end{equation*}
con igualdad si y sólo si
\begin{equation*}
    \gamma(t)=c_0+c_1e^{it},
    \qquad |c_1|=\frac{L}{2\pi}.
\end{equation*}
Esta curva es un círculo con centro en $c_0$ y radio $|c_1|$.
\hfill $\square$
\begin{tcolorbox}[title=\textcolor{white}{Problema 17}]
\textit{¿Dónde usaste que la curva fuera simple? Observa que si defines}
\begin{equation*}
    A=\frac12\int_{-\pi}^{\pi}\bigl(x(t)y'(t)-x'(t)y(t)\bigr)\,dt,
\end{equation*}
\textit{la desigualdad sigue siendo válida y no necesitamos suponer que la curva es simple. Usa esta observación para probar el siguiente resultado.}

\medskip
\textbf{Teorema 3 (Desigualdad de Wirtinger).} \textit{Sea $f:\R\to\R$ una función periódica de clase $C^2$ de periodo $2\pi$, tal que}
\begin{equation*}
    \int_{-\pi}^{\pi} f(t)\,dt=0.
\end{equation*}
\textit{Entonces}
\begin{equation*}
    \int_{-\pi}^{\pi} f^2(t)\,dt \leq \int_{-\pi}^{\pi} \bigl(f'(t)\bigr)^2\,dt,
\end{equation*}
\textit{con igualdad si y sólo si}
\begin{equation*}
    f(t)=a\cos t+b\sen t.
\end{equation*}
\end{tcolorbox}
\textbf{Demostración.}\,\
Sea
\begin{equation*}
    f(t)\sim \sum_{n\in\Z}c_ne^{int}
\end{equation*}
la serie de Fourier de $f$. Como $f$ es de clase $C^2$, podemos derivar su serie de Fourier y escribir
\begin{equation*}
    f'(t)\sim i\sum_{n\in\Z}n c_ne^{int}.
\end{equation*}
Además, como $f$ es real,
\begin{equation*}
    \overline{f(t)}=f(t)
    \qquad\text{y}\qquad
    \overline{f'(t)}=f'(t).
\end{equation*}
Por tanto,
\begin{equation*}
    \overline{f(t)}\sim \sum_{k\in\Z}\overline{c_k}e^{-ikt},
\end{equation*}
y
\begin{equation*}
    \overline{f'(t)}
    \sim -i\sum_{k\in\Z}k\overline{c_k}e^{-ikt}.
\end{equation*}

Primero calculamos la integral de $\bigl(f'(t)\bigr)^2$. Como $f'$ es real,
\begin{equation*}
    \bigl(f'(t)\bigr)^2=f'(t)\overline{f'(t)}.
\end{equation*}
Así,
\begin{align*}
    \int_{-\pi}^{\pi}\bigl(f'(t)\bigr)^2\,dt
    &=\int_{-\pi}^{\pi}
    \left(i\sum_{n\in\Z}n c_ne^{int}\right)
    \left(-i\sum_{k\in\Z}k\overline{c_k}e^{-ikt}\right)\,dt \\
    &=\sum_{n\in\Z}\sum_{k\in\Z}nk c_n\overline{c_k}
    \int_{-\pi}^{\pi}e^{i(n-k)t}\,dt.
\end{align*}
Por la ortogonalidad del Problema 3(a),
\begin{equation*}
    \int_{-\pi}^{\pi}e^{i(n-k)t}\,dt=
    \begin{cases}
        2\pi, & n=k,\\
        0, & n\neq k.
    \end{cases}
\end{equation*}
Entonces
\begin{equation*}
    \int_{-\pi}^{\pi}\bigl(f'(t)\bigr)^2\,dt
    =2\pi\sum_{n\in\Z}n^2|c_n|^2.
\end{equation*}

Ahora calculamos la integral de $f^2(t)$. Como $f$ es real,
\begin{equation*}
    f^2(t)=f(t)\overline{f(t)}.
\end{equation*}
Entonces
\begin{align*}
    \int_{-\pi}^{\pi}f^2(t)\,dt
    &=\int_{-\pi}^{\pi}
    \left(\sum_{n\in\Z}c_ne^{int}\right)
    \left(\sum_{k\in\Z}\overline{c_k}e^{-ikt}\right)\,dt \\
    &=\sum_{n\in\Z}\sum_{k\in\Z}c_n\overline{c_k}
    \int_{-\pi}^{\pi}e^{i(n-k)t}\,dt \\
    &=2\pi\sum_{n\in\Z}|c_n|^2.
\end{align*}

Como
\begin{equation*}
    \int_{-\pi}^{\pi}f(t)\,dt=0,
\end{equation*}
tenemos
\begin{equation*}
    c_0=\frac{1}{2\pi}\int_{-\pi}^{\pi}f(t)\,dt=0.
\end{equation*}
Por tanto,
\begin{equation*}
    \sum_{n\in\Z}|c_n|^2
    =\sum_{\substack{n\in\Z\\ n\neq0}}|c_n|^2.
\end{equation*}
Como si $n\neq0$ entonces $n^2\geq1$, se sigue que
\begin{equation*}
    n^2|c_n|^2\geq |c_n|^2.
\end{equation*}
Así,
\begin{equation*}
    \sum_{n\in\Z}n^2|c_n|^2
    \geq
    \sum_{n\in\Z}|c_n|^2.
\end{equation*}
Multiplicando por $2\pi$, obtenemos
\begin{equation*}
    2\pi\sum_{n\in\Z}n^2|c_n|^2
    \geq
    2\pi\sum_{n\in\Z}|c_n|^2.
\end{equation*}
Es decir,
\begin{equation*}
    \int_{-\pi}^{\pi}\bigl(f'(t)\bigr)^2\,dt
    \geq
    \int_{-\pi}^{\pi}f^2(t)\,dt.
\end{equation*}
Por lo tanto,
\begin{equation*}
    \int_{-\pi}^{\pi} f^2(t)\,dt
    \leq
    \int_{-\pi}^{\pi}\bigl(f'(t)\bigr)^2\,dt.
\end{equation*}

Ahora veamos el caso de igualdad. La igualdad ocurre si y sólo si
\begin{equation*}
    \sum_{n\in\Z}\bigl(n^2-1\bigr)|c_n|^2=0.
\end{equation*}
Como cada término es no negativo, esto sucede si y sólo si
\begin{equation*}
    c_n=0
    \qquad\text{para todo } n\neq -1,0,1.
\end{equation*}
Pero además $c_0=0$, así que
\begin{equation*}
    f(t)=c_1e^{it}+c_{-1}e^{-it}.
\end{equation*}
Como $f$ es real, se cumple $c_{-1}=\overline{c_1}$. Si escribimos
\begin{equation*}
    c_1=\frac{a-ib}{2},
\end{equation*}
entonces
\begin{align*}
    f(t)
    &=\frac{a-ib}{2}e^{it}+\frac{a+ib}{2}e^{-it}\\
    &=a\cos t+b\sen t.
\end{align*}
Recíprocamente, si
\begin{equation*}
    f(t)=a\cos t+b\sen t,
\end{equation*}
entonces sus únicos coeficientes de Fourier no nulos están en $n=-1$ y $n=1$, y por tanto
\begin{equation*}
    \sum_{n\in\Z}\bigl(n^2-1\bigr)|c_n|^2=0.
\end{equation*}
Así, en ese caso se tiene la igualdad.
\hfill $\square$
\begin{tcolorbox}[title=\textcolor{white}{Problema 18}]
\textit{Prueba que para una curva $C^2$, $\gamma(t)=(x(t),y(t))$ uno tiene que:
\begin{equation*}
    \int_{-\pi}^{\pi}[x'(t)+y'(t)]^2dt+\int_{-\pi}^{\pi}((y'(t))^2-y^2(t))dt
\end{equation*}}
\end{tcolorbox}
\textbf{Demostraci\'on}:Del problema anterior se sique que $\int_{-\pi}^{\pi}((y'(t))^2-y^2(t))dt>0$,pues, entonces
\begin{equation*}
     \int_{-\pi}^{\pi}[x'(t)+y'(t)]^2dt+\int_{-\pi}^{\pi}((y'(t))^2-y^2(t))dt>0
\end{equation*}
luego:
\begin{equation*}
     \int_{-\pi}^{\pi}[x'(t)+y'(t)]^2dt+\int_{-\pi}^{\pi}((y'(t))^2-y^2(t))dt=\int_{-\pi}^{\pi}(y'(t))^2+(x'(t))^2dt+2\int_{-\pi}^{\pi}x'(t)y(t)=\int_{-\pi}^{\pi}\frac{L^2}{4\pi^2}-2A=\frac{L^2}{2\pi}-2A\geq 0
\end{equation*}  
\hfill $\square$
\begin{tcolorbox}[title=\textcolor{white}{Problema 19}]
\textit{Sea $\gamma:[a,b]\to\R^2$, $\gamma(t)=(x(t),y(t))$, una curva cerrada de clase $C^1$. Muestre que}
\begin{equation*}
    \frac12\int_a^b\bigl(x(t)y'(t)-x'(t)y(t)\bigr)\,dt
    =\int_a^b x(t)y'(t)\,dt
    =-\int_a^b x'(t)y(t)\,dt.
\end{equation*}
\end{tcolorbox}

\textbf{Demostración.}\,\
Partimos del lado izquierdo:
\begin{align*}
    \frac12\int_a^b\bigl(x(t)y'(t)-x'(t)y(t)\bigr)\dt
    &=\frac12\int_a^b x(t)y'(t)\dt
      -\frac12\int_a^b x'(t)y(t)\dt.
\end{align*}
Integramos por partes el segundo término. Como
\begin{equation*}
    \int_a^b x'(t)y(t)\dt
    =\left[x(t)y(t)\right]_a^b-\int_a^b x(t)y'(t)\dt,
\end{equation*}
entonces
\begin{align*}
    \frac12\int_a^b\bigl(x(t)y'(t)-x'(t)y(t)\bigr)\dt
    &=\frac12\int_a^b x(t)y'(t)\dt
    -\frac12\left(\left[x(t)y(t)\right]_a^b-\int_a^b x(t)y'(t)\dt\right) \\
    &=\int_a^b x(t)y'(t)\dt
    -\frac12\left[x(t)y(t)\right]_a^b.
\end{align*}
Como $\gamma$ es cerrada, se cumple
\begin{equation*}
    \gamma(a)=\gamma(b).
\end{equation*}
En particular,
\begin{equation*}
    x(a)=x(b)
    \qquad\text{y}\qquad
    y(a)=y(b).
\end{equation*}
Por tanto,
\begin{equation*}
    \left[x(t)y(t)\right]_a^b=x(b)y(b)-x(a)y(a)=0.
\end{equation*}
Finalmente,
\begin{equation*}
    \frac12\int_a^b\bigl(x(t)y'(t)-x'(t)y(t)\bigr)\dt
    =\int_a^b x(t)y'(t)\dt.
\end{equation*}
\hfill $\square$
\begin{tcolorbox}[title=\textcolor{white}{Problema 20}]
\textit{Define $\gamma^-:[a,b]\to\R^2$ como}
\begin{equation*}
    \gamma^-(t)=\gamma(a+b-t).
\end{equation*}
\textit{La imagen de $\gamma^-$ es la misma que la de $\gamma$, pero recorrida en la dirección opuesta. Muestra que, para $\gamma^-$, la integral del Teorema 4 cambia de signo.}
\end{tcolorbox}

\textbf{Demostración.}\,
Sea
\begin{equation*}
    \gamma(t)=(x(t),y(t)).
\end{equation*}
Entonces
\begin{equation*}
    \gamma^-(t)=\gamma(a+b-t)=\bigl(x(a+b-t),y(a+b-t)\bigr).
\end{equation*}
Denotemos
\begin{equation*}
    x^-(t)=x(a+b-t),
    \qquad
    y^-(t)=y(a+b-t).
\end{equation*}
Por la regla de la cadena,
\begin{equation*}
    (x^-)'(t)=-x'(a+b-t),
    \qquad
    (y^-)'(t)=-y'(a+b-t).
\end{equation*}

La integral del Teorema 4 para $\gamma^-$ es
\begin{equation*}
    I^-=\frac12\int_a^b\bigl(x^-(t)(y^-)'(t)-(x^-)'(t)y^-(t)\bigr)\,dt.
\end{equation*}
Sustituyendo las expresiones anteriores,
\begin{align*}
    I^-
    &=\frac12\int_a^b\Bigl(x(a+b-t)\bigl[-y'(a+b-t)\bigr]
    -\bigl[-x'(a+b-t)\bigr]y(a+b-t)\Bigr)\,dt \\
    &=\frac12\int_a^b\Bigl(-x(a+b-t)y'(a+b-t)
    +x'(a+b-t)y(a+b-t)\Bigr)\,dt.
\end{align*}
Hacemos el cambio de variable
\begin{equation*}
    u=a+b-t,
    \qquad du=-dt.
\end{equation*}
Cuando $t=a$, se tiene $u=b$; cuando $t=b$, se tiene $u=a$. Entonces
\begin{align*}
    I^-
    &=\frac12\int_b^a\bigl(-x(u)y'(u)+x'(u)y(u)\bigr)(-du) \\
    &=\frac12\int_a^b\bigl(-x(u)y'(u)+x'(u)y(u)\bigr)\,du \\
    &=-\frac12\int_a^b\bigl(x(u)y'(u)-x'(u)y(u)\bigr)\,du.
\end{align*}
Por tanto, si
\begin{equation*}
    I=\frac12\int_a^b\bigl(x(t)y'(t)-x'(t)y(t)\bigr)\,dt,
\end{equation*}
entonces
\begin{equation*}
    I^-=-I.
\end{equation*}
Así, para $\gamma^-$, la integral del Teorema 4 cambia de signo.
\hfill $\square$


\newpage

\begin{tcolorbox}[title=\textcolor{white}{Problema 21}]
\textit{Supón que la curva descrita por $\gamma$ es como en la Figura 1, es decir,}
\begin{equation*}
    \gamma(a)=(0,0),
    \qquad
    \gamma(T)=(x(T),0)
\end{equation*}
\textit{para algún $T\in(a,b)$; que la imagen de $\gamma$ en el intervalo $[a,T]$ es la gráfica de una función de clase $C^1$ $g$, y que la imagen de $\gamma$ en el intervalo $[T,b]$ es la gráfica de otra función de clase $C^1$ $f$, con}
\begin{equation*}
    f(x)\geq g(x),
    \qquad
    x\in[0,x(T)].
\end{equation*}
\textit{Argumenta que}
\begin{equation*}
    A=\int_0^{x(T)}\bigl(f(x)-g(x)\bigr)\,dx
\end{equation*}
\textit{y prueba, usando cambio de variable, que}
\begin{equation*}
    \int_0^{x(T)}g(x)\,dx
    =
    \int_a^T y(t)x'(t)\,dt
\end{equation*}
\textit{y}
\begin{equation*}
    \int_0^{x(T)}f(x)\,dx
    =
    -\int_T^b y(t)x'(t)\,dt.
\end{equation*}
\textit{Concluye que el Teorema 4 es cierto en este caso.}
\end{tcolorbox}

\textbf{Demostración.}\,\
Como la región encerrada por la curva queda entre la gráfica superior $y=f(x)$ y la gráfica inferior $y=g(x)$, con
\begin{equation*}
    0\leq x\leq x(T),
\end{equation*}
su área está dada por
\begin{equation*}
    A=\int_0^{x(T)}\bigl(f(x)-g(x)\bigr)\,dx.
\end{equation*}

En el tramo $[a,T]$, la curva coincide con la gráfica de $g$. Por tanto,
\begin{equation*}
    y(t)=g(x(t)).
\end{equation*}
Usando el cambio de variable
\begin{equation*}
    x=x(t),
    \qquad
    dx=x'(t)\,dt,
\end{equation*}
obtenemos
\begin{align*}
    \int_0^{x(T)}g(x)\,dx
    &=\int_a^T g(x(t))x'(t)\,dt\\
    &=\int_a^T y(t)x'(t)\,dt.
\end{align*}

Ahora, en el tramo $[T,b]$, la curva coincide con la gráfica de $f$, pero se recorre en sentido contrario, desde $x(T)$ hasta $0$. En este tramo se cumple
\begin{equation*}
    y(t)=f(x(t)).
\end{equation*}
Así,
\begin{align*}
    \int_0^{x(T)}f(x)\,dx
    &=-\int_{x(T)}^0 f(x)\,dx\\
    &=-\int_T^b f(x(t))x'(t)\,dt\\
    &=-\int_T^b y(t)x'(t)\,dt.
\end{align*}

Entonces, sustituyendo las dos igualdades anteriores en la expresión del área,
\begin{align*}
    A
      &=\int_0^{x(T)}f(x)\,dx-\int_0^{x(T)}g(x)\,dx\\
    &=-\int_T^b y(t)x'(t)\,dt-\int_a^T y(t)x'(t)\,dt\\
    &=-\int_a^b y(t)x'(t)\,dt.
\end{align*}
Por el Problema 19,
\begin{equation*}
    -\int_a^b y(t)x'(t)\,dt
    =
    \frac12\int_a^b\bigl(x(t)y'(t)-x'(t)y(t)\bigr)\,dt.
\end{equation*}
Por lo tanto,
\begin{equation*}
    A=
    \frac12\int_a^b\bigl(x(t)y'(t)-x'(t)y(t)\bigr)\,dt.
\end{equation*}
Así, el Teorema 4 queda probado en este caso.
\hfill $\square$
\begin{tcolorbox}[title=\textcolor{white}{Problema 22}]
\textit{Demuestra que lo anterior se puede generalizar al caso en que la region se subdivide en un n\'umero finito de partes}
\end{tcolorbox}
\textbf{Demostración:} Simplemente usamos la propiedad de linealidad de la integral, pues las regiones son finitas, en cada una se vale el teorema anterior 
\hfill $\square$
\end{document}